% !PRINTED 633
\weberpage{633}
%<<CONT
Wurzel von $X = 0$ setzt, also die Substitution $(\varepsilon, \varepsilon^h)$ für $h = 1, 2 \ldots n-1$ ausführt. Alle diese Körper sind wieder in $\Omega(\varepsilon)$ enthalten; sie sind also mit einander identisch, d.\,h. $\Omega(\varepsilon)$ ist ein \gesperrt{Normalkörper}. Wenn eine Grösse $\Phi(\varepsilon)$ die Eigenschaft hat, ungeändert zu bleiben, wenn $\varepsilon$ durch alle $\varepsilon^h$ ersetzt wird, so ist sie selbst in $\Omega$ enthalten. Dies folgt schon aus den allgemeinen Sätzen des \S.\,144. Wir sehen aber die Richtigkeit sofort ein, wenn wir unter der Voraussetzung $\Phi(\varepsilon^h) = \Phi(\varepsilon)$
\[
\Phi(\varepsilon) = \frac{1}{n-1}\sum_{1,\,n-1}^{h}\Phi(\varepsilon^h) = c_0 - \frac{c_1 + c_2 + \cdots + c_{n-2}}{n-1}
\]
setzen, woraus dann nach 2. zu schliessen ist, dass $c_1, c_2 \ldots c_{n-2}$ gleich Null sind. Also:

\gesperrt{3. Wenn eine Grösse in $\Omega(\varepsilon)$ die sämmtlichen Substitutionen $(\varepsilon, \varepsilon^h)$ gestattet, so ist sie in $\Omega$ enthalten.}

\begin{center}
\S.\,184.

\gesperrt{Sätze über die Resolventen.}
\end{center}

Wir betrachten jetzt die \gesperrt{Lagrange}'schen Resolventen $(\varepsilon, x)$ unter der Voraussetzung, dass $\varepsilon$ eine \gesperrt{imaginäre} $n^{\text{te}}$ Einheitswurzel und die $x$ unabhängige Variable sind, und wenden darauf die Sätze des vorigen Paragraphen an.

Wir untersuchen zunächst den Einfluss der linearen Permutationen auf $(\varepsilon, x)$; dazu genügt es, wenn wir den Einfluss der beiden erzeugenden Substitutionen
\weq{1}{s = (h, h+1), \quad t = (h, gh)}
feststellen, worin das nach dem Modul $n$ genommene $h$ die Indices von $x$ durchläuft und $g$ eine primitive Wurzel der Primzahl $n$ bedeutet (\S.\,180).

Wenden wir aber auf
\weq{2}{(\varepsilon, x) = x_0 + \varepsilon x_1 + \varepsilon^2 x_2 + \cdots + \varepsilon^{n-1}x_{n-1}}
die Substitution $s$ an, so geht diese Function über in
\[
x_1 + \varepsilon x_2 + \varepsilon^2 x_3 + \cdots + \varepsilon^{n-1}x_0 = \varepsilon^{-1}(\varepsilon, x).
\]
Der Einfluss der Substitution $s$ ist also, dass
\[
(\varepsilon, x) \quad \text{in} \quad \varepsilon^{-1}(\varepsilon, x)
\]
übergeht.
